\begin{abstract}
Distributed trace anomaly scoring commonly represents each trace as a
Service Trace Vector (STV), where features correspond to service call
paths and values encode latency. However, raw STV-based scoring can be
strongly influenced by trace structural complexity. In this paper, we
show empirically on three TrainTicket configurations of
\texttt{ts-order-service} that baseline STV~+~Isolation Forest scores
correlate strongly with trace size, with correlations of approximately
\(r=0.85\)--\(0.98\) for span and service-count based complexity
measures. We propose WR-STV, a lightweight STV-based anomaly scoring
method that replaces raw feature scoring with clipped standardized
residuals, applies log-frequency reliability weighting, and aggregates
the top-\(k\) weighted residual contributors into a trace-level score.
On held-out non-injected traces, WR-STV reduces the measured
score-complexity correlation from approximately \(r=0.87\) to
\(r=0.30\). Under controlled STV-visible latency injection, WR-STV
improves within-dataset ROC-AUC from \(0.605\) to \(0.814\), with
random-visible injection used to test whether the improvement depends
on perturbing only high-duration spans. Additional baseline comparisons
indicate that residual normalization accounts for most of the gain,
while reliability weighting provides a smaller, metric-dependent
refinement. WR-STV also provides interpretable service-path-level
evidence and lightweight per-trace scoring. The evaluation is scoped to
synthetic latency perturbations on three configurations of one
benchmark service.
\end{abstract}